% !TEX root = ../main.tex

% 预习报告	
\setcounter{section}{0}
\section{APL1-3 Quantum Key Distribution \\ Preview Report}


%---------------------------------------------------------------------
% 实验目的
\subsection{Purpose}
\begin{enumerate}
	\item 掌握控制和测量光子的偏振；
	\item 掌握单光子的标定；
	\item 掌握单光子的探测及相应探测器效率的测量；
	\item 掌握 BB84 量子密钥分发过程的数据处理。
\end{enumerate}


%---------------------------------------------------------------------
% 仪器用具
\subsection{Instruments \& Equipment}
\begin{table}[htbp]
	\centering
	\renewcommand\arraystretch{1.6}
	% \setlength{\tabcolsep}{10mm}
	\caption{偏振测量实验}
	\begin{tabular}{p{0.05\textwidth}|p{0.20\textwidth}|p{0.05\textwidth}|p{0.5\textwidth}}
		\hline
		编号& 仪器用具名称 & 数量 &  主要参数（型号，测量范围，测量精度等） \\
		\hline
		1	&	准直激光器 	& 1 & 波长：404nm，最大功率：150mW\\
		
		2	&	偏振分光棱镜 	& 2 & 波长：404nm，消光比＞500	 \\
		
		3	&	半波片 & 2 &	波长：404nm，零级 	\\
		
		4	&	小型磁性底座	&  & MB105	\\
		
		5	&	PH 系列杆架	& 6 & PH102	\\
		
		6	&	SP 系列接杆	& 6 & SP104 \\
		
		7	&	激光器镜架	& 6 & OM311	\\
		
		8	&	精密棱镜台	& 2 & PPM101	\\
		
		9	&	偏光镜架	& 2 & PM101	\\
		
		10	&	可见光功率计& 2 & PM100、S120VC	\\
		
		11	&	直流稳压电源& 1 & GPD-3303D	\\
		\hline
	\end{tabular}
\end{table}

\begin{table}[htbp]
	\centering
	\renewcommand\arraystretch{1.6}
	% \setlength{\tabcolsep}{10mm}
	\caption{单光子标定的用具}
	\begin{tabular}{p{0.05\textwidth}|p{0.20\textwidth}|p{0.05\textwidth}|p{0.5\textwidth}}
		\hline
		编号& 仪器用具名称 & 数量 &  主要参数（型号，测量范围，测量精度等） \\
		\hline
		1	&	密钥分发系统 	& 1 & 波长：404nm	\\
		
		2	&	可见光功率计 	& 1 & PM100、S120VC	 \\
		\hline
	\end{tabular}
\end{table}

\begin{table}[htbp]
	\centering
	\renewcommand\arraystretch{1.6}
	% \setlength{\tabcolsep}{10mm}
	\caption{单光子的探测及相应探测器探测效率测量的用具}
	\begin{tabular}{p{0.05\textwidth}|p{0.20\textwidth}|p{0.05\textwidth}|p{0.5\textwidth}}
		\hline
		编号& 仪器用具名称 & 数量 &  主要参数（型号，测量范围，测量精度等） \\
		\hline
		1	&	反射镜 	& 1 & 波长：404nm，45 度入射	\\
		
		2	&	滤波片 	& 1 & 波长：405nm，带宽：3nm	 \\
		
		3	&	光纤准直器 & 1 &	F671FC-405 	\\
		
		4	&	反射镜折叠架	& 1 & OM402	\\
		
		5	&	透镜固定架	& 1 & LH102	\\
		
		6	&	光纤耦合架	& 1 & PFC201 \\
		
		7	&	小型磁性底座	& 3 & MB105	\\
		
		8	&	PH 系列杆架	& 3 & PH102	\\
		
		9	&	SP 系列接杆	& 1 & SP104	\\
		
		10	&	SP 系列接杆	& 2	& SP134	\\
		
		11	&	可见光功率计	& 1 & PM100、S120VC	\\
		
		12	&	直流稳压电源& 1 & GPD-3303D	\\
		\hline
	\end{tabular}
\end{table}


\begin{table}[htbp]
	\centering
	\renewcommand\arraystretch{1.6}
	% \setlength{\tabcolsep}{10mm}
	\caption{密钥分发过程数据处理的用具}
	\begin{tabular}{p{0.05\textwidth}|p{0.20\textwidth}|p{0.05\textwidth}|p{0.5\textwidth}}
		\hline
		编号& 仪器用具名称 & 数量 &  主要参数（型号，测量范围，测量精度等） \\
		\hline
		1	&	密钥分发系统 	& 1 & 波长：404nm	\\
		\hline
	\end{tabular}
\end{table}


\clearpage
%---------------------------------------------------------------------
% 原理概述
\subsection{Principle}
\begin{enumerate}
	\item 量子密钥分发及 BB84 协议
	\begin{itemize}
		\item 概述
		
		BB84协议是由Charles Bennett和Gilles Brassard在1984年提出的一种量子秘钥分发（QKD）协议。其核心思想是利用量子力学的原理来安全地分配密钥，确保通信的安全性。
		
		BB84协议通过量子比特（qubit）来传输信息，使用光子的极化态作为信息载体。通信双方通常称为Alice（发送者）和Bob（接收者）。Alice随机选择一个比特（0或1），并使用两个不同的基（基1和基2）进行编码。基1可以是水平（|0>）和垂直（|1>）极化，基2可以是对角线（|+>）和反对角线（|−>）极化。Alice为每个比特随机选择基进行编码。例如，如果她选择发送比特0并使用基1，她会发送水平极化的光子；如果选择基2，她会发送对角线极化的光子。Alice将准备好的光子通过量子信道发送给Bob。这个过程是量子通信的核心，任何对光子状态的测量都会引入干扰。Bob收到光子后，他也随机选择一个基进行测量。由于他对基的选择是随机的，可能会与Alice的基相同，也可能不同。如果Bob的基与Alice的基相同，他将准确地测量到Alice发送的比特。如果基不同，测量结果将是随机的。Alice和Bob通过经典信道（不安全）交换他们选择的基的信息。双方保留在同一基下进行的测量结果，而丢弃不匹配的测量结果。 Alice和Bob现在都有了一串比特，它们对应于相同基下的测量结果。这些比特可以作为共享的秘密密钥。
		
		BB84协议的安全性基于量子力学的原理，如量子不可克隆定理和测量扰动原则。量子不可克隆定理表明，无法精确复制一个未知的量子状态，因此Eve无法获取完整的密钥信息。此外，测量光子的状态会改变其状态，导致Alice和Bob的测量结果不一致，从而能检测到窃听。
		
		\item 详解
		
		\textbf{以下内容来自CSDN博文“BB84协议:量子秘钥分发”：}
		\begin{itemize}
			\item 基本通讯
			
			首先，BB84协议要求通信双方拥有一条量子信道的同时，还拥有一条传统信道。然后在量子信道上，信息发送方需要4种状态的量子作为信使，信息接收方需要2种量子接收器。两组量子信使需要分别使用两种不同的接收器来接收。“横竖”组对应“+”形接收器，“撇捺”组对应“x”形接收器。我们规定每组量子中有一种姿态代表0，另一种代表 1。例如，我们规定，接收到“横”代表0，那么，接收到“竖”代表 1。同理，规定接收到“撇”代表0，接收到“捺”代表 1。
			
			现在如果我们挑选正确的接收器，那么我们能正确地接收到信息发送者想发给我们的信息。例如，我们如果用“+”形接收器来接收姿态为“横”的量子，那么我们得到的信号就是“横”，那么发送者想传递的“0”这个信号就会被我们正确地接收到。如果我们没有用正确的接收器会发生什么呢？如果我们用“x”形接收器接收姿态为“横”的量子会发生什么呢？量子信使会撞在接收器的入口上，导致接收器接收不到信使吗？不一定。姿态为“横”的量子经过“x”形接收器时，有可能会改变自己的姿态，变成能通过接收器的姿态，也就是姿态为“横”的量子经过“x”形接收器时，姿态可能变为“撇”，也可能变为“捺”，那么传达的信息有可能还是发送者想传达的0，但也可能变成了相反的 1。错误地使用“+”形接收器接收“撇捺”组量子信使同理，信息接收者无法分辨原本信息的发送者想传递的究竟是0还是 1。
			
			按密码学传统，我们把信息发送者叫做Alice，信息接收者叫做Bob。现在他们要用BB84协议，为他们的通信确定一个密钥。注意在量子信道上的通讯其实到步骤3就结束了，接下来的步骤4-7都是在传统通信信道上进行的。我们认为传统信道意味着攻击者可能窃听并修改信道中会话内容。
			\begin{enumerate}
				\item 第一步，Alice选出一串二进制串（01101100001001…），命名为S。
				\item 第二步，Alice逐位地、随机地为S的每一位指定量子信使的组。例如S的某一位是 1，Alice为这一位指定的量子信使组是“撇捺”组。那么稍后这一位上的 1就会用姿态为“捺”的量子信使发射出去，而不用姿态为“横”的信使。
				\item 第三步，Bob随机地用一系列的接收器接收Alice发来的量子信使，得到一系列0、1二进制串。由于Bob并不知道Alice发送的量子信使种类是怎么排列的，可以想见的是，Bob会在很多位上使用错误的量子接收器接收量子信使。例如发送的二进制串的某一位上是 1，Alice用的是“横竖”组的竖姿态量子信使发送这一位的 1，但是Bob可能选择了“x”形接收器来接收这一位信息。结果可能是Bob得到了错误的信号0，（当然也可能是碰巧还是得到了正确的 1，因为“x”形接收器接收竖姿态量子，可能得到撇姿态量子，也可能得到捺姿态量子。）
				\item 第四步，Bob通过传统信道中告诉Alice自己的测量方法，即自己用哪种接收器来接收S的每一位信息。例如，Bob会告诉Alice：我第一位用的“+”形接收器、第二位用的“x”形、第三位用的“x”形…Bob在这一步只告诉Alice自己是怎么测量的，没有告诉Alice自己测量的结果。
				\item 第五步，Alice收到Bob的测量方法后，对比自己的发送量子信使的情况，告诉Bob他在哪些位上碰巧用了正确的接收器。我们把正确测量的位的集合称为C，Alice把C发送给Bob。例如，Alice发现第2位自己发射的是撇姿态量子，而Bob正好是用“x”形接收器接收的量子，那么Bob第二位的测量是正确的，则第2位属于C。然后，Alice发现自己第3位发送的是竖姿态量子，而Bob在这一位上使用的是“x”形接收器，那么Bob第三位的测量方式就是错误的，则第3位不属于C。
				\item 第六步，Bob收到Alice的信息，知道了自己哪些位的测量方式是正确的，即知道了集合C。他在自己测量正确的那些位中，随机挑选出一部分位，把测量到的结果发送给Alice。我们把Bob从C中挑出的那一部分叫集合D。例如Bob从Alice的信息中得知自己第6、8、17、45、46、56位的测量方式是正确的（集合C），那么他随机挑出几位，例如，挑出第6、17、45位，这就是集合D。Bob将属于D的这几位上自己测到的结果（是0还是 1）发送给Alice：“我在第6、17、45位分别测到了 1、1、0”。
				\item 第七步，Alice收到信息后核对自己在6、17、45位上发送的信息是不是如Bob所说的 1、1、0。如果是，那么Alice发送信息给Bob示意没问题，那么双方会在Bob一开始测量正确的位（集合C）中，剔除第六步Bob发回给Alice的那些位（集合D），剩下的位（D在C中的补集）作为双方通信的密钥。按上述例子，Bob在第6、8、17、45、46、56位的测量方式是正确的，他又把第6、17、45位发回给Alice，那么用作通信密钥的就是剩下的8、46、56位的值。例如Alice最开始发送的二进制字符串S中第8、46、56位值分别为 1、0、1，由于这几位Bob用了正确的接收器，所以Bob接收到的也是 1、0、1，双方都会以“101”作为通讯密钥，密钥协商的目的达成。
			\end{enumerate}
			
			\item 为什么这样做能防止通讯被窃听?
			\begin{itemize}
				\item 在传统信道上的窃听
				
				按照密码学传统，我们称窃听者为Eve。如果Eve在传统信道上窃听，她会在上述第四步听到Bob的测量方式，但是Bob没有发送测量结果。Eve还会听到Bob一部分正确测量位以及测出的结果（上述第六步中，Eve会听到集合D以及对应的测量值），但是最后Alice与Bob选择通信密钥时会完全丢弃集合D，即完全丢弃在传统信道交流过的那些位。到最后，Eve只知道Alice与Bob使用了S的哪些位用作通讯密钥，但是Eve不知道这些位上的具体数值是多少。
				
				\item 在量子信道上的窃听
				
				\begin{figure}[htbp]
					\centering
					\includegraphics[width=0.6\textwidth]{APL1_3_pre_Eve.png}
					\caption{在量子信道上的窃听}
				\end{figure}
				
				Eve遇到上面的困难，心一横，决定去窃听量子信道上Alice一开始传给Bob的二进制串S。但是首先Eve就遇到一个问题：她不知道该选用什么接收器序列来接收量子信使（连Bob都不知道）。Eve又心一横，用了一组随机的量子接收器序列接收量子信使。这样一来Eve必然使用了很多不合适的接收器，这样也就导致了一些量子信使的状态被改变，这些改变会在Alice-Bob后续的交流中被察觉。
				
				例如，S在某一位上的值是 1，Alice用竖姿态量子信使发送了这一位信息，但是Eve错用“x”形接收器接收这一位信息，这样此量子可能变成了撇量子（或捺量子），但Bob在这一位上恰好猜对了，他用“+”形接收器接收经过Eve窃听过的量子信使（撇或捺量子信使），结果量子信使变成了横姿态量子信使（当然也有接收成竖姿态量子，信息碰巧正确地被传递了，这时这一位不会引起Alice-Bob的注意），原本是 1的值Bob用合适的接收器接收却得到了0。然后在上述协议的第六步，选集合C的时候，如果Bob正好把这一位选中发给了Alice，那么Alice就会发现Bob在这一位上明明使用了正确的接收器却没有得到正确的值。这种情况出现频率超过某个阈值，Alice就会察觉有人在窃听量子信道。
			\end{itemize}
			
			所以其实BB84协议只能保证通讯被窃听的时候，通讯双方能及时发现，之后要么切换信道，要么终止通信。
				
		\end{itemize}	
	\end{itemize}
	
	\item 光的偏振的测量和控制
	
	在量子密钥分发实验中，偏振分光棱镜、半波片和探测器的联合使用能够有效地制备、控制和测量光的偏振态。特别是，半波片的旋转可以精确控制偏振方向，而PBS的功率探测可以用于验证光子的偏振态。
	\begin{itemize}
		\item 偏振分光棱镜 (PBS)
		
		偏振分光棱镜是一种根据光的偏振态对光束进行分离的光学器件。当一束光垂直入射到偏振分光棱镜的表面时，光束中的水平偏振分量将透过棱镜，而垂直偏振分量则会被反射。这意味着透射端口输出的是水平偏振光（H态），而反射端口输出的是垂直偏振光（V态）。在经典光学中，这种分光效果可以解释为光波的电场矢量在水平和垂直方向上的分量不同，因此它们被不同地处理。
		
		\begin{figure}[htbp]
			\centering
			\includegraphics[width=0.6\textwidth]{APL1_3_pre_PBS_and_waveplate.png}
			\caption{PBS和波片}
		\end{figure}
		
		在量子密钥分发（QKD）的实验中，光子的偏振态可以描述为量子叠加态，例如，水平偏振态 \( |H\rangle \) 和垂直偏振态 \( |V\rangle \) 的线性叠加态。单个光子的偏振态在宏观表现上对应于一束光的整体偏振状态。例如，光子的偏振态 \( |D\rangle \) 可以用来表示与水平面成 45° 夹角的线偏振光。
		
		\item 起偏与偏振态的制备
		
		在实验中，起偏器件（如偏振分光棱镜）可以用于产生纯粹的偏振态。具体来说，光束首先通过PBS的透射端口，输出水平偏振光 \( |H\rangle \)。为了产生其他角度的线偏振光，实验中通常使用半波片。半波片是一种双折射晶体，它可以通过改变光在两个正交偏振方向上的传播速度，从而旋转光的偏振态。如果将半波片旋转角度为 \( \theta \)，则入射的偏振光的偏振方向将旋转 \( 2\theta \)。例如，要制备偏振方向为 45° 的线偏振光 \( |D\rangle \)，半波片需要旋转 22.5°，这样可以将水平偏振光 \( |H\rangle \) 转变为 \( |D\rangle \)。
		
		\item 偏振态的测量
		
		测量45°的线偏振光通常需要结合PBS和探测器。在实验中，通过PBS可以分别探测到水平和垂直偏振光的分量。例如，将 45° 偏振光通过一个PBS，理论上在透射端和反射端探测到的功率比应为1:1。换句话说，水平偏振光 \( |H\rangle \) 和垂直偏振光 \( |V\rangle \) 各占50\%的强度。这是由于 45° 偏振光 \( |D\rangle \) 可以表示为两个正交偏振态的等量叠加：
		\[
		|D\rangle = \frac{1}{\sqrt{2}} (|H\rangle + |V\rangle)
		\]
		因此，探测到的功率会在两个端口上均分。
		
		然而，为了进一步验证45°偏振光的存在，可以先通过一个22.5°旋转的半波片，这样能够将45°偏振光 \( |D\rangle \) 转化为纯水平偏振光 \( |H\rangle \)。当这种偏振光再经过PBS时，透射端应当检测到全部的功率，而反射端则检测不到任何功率。也就是说，反射端的功率为零。这一实验现象能够更有效地验证45°偏振光的存在，因为在这种情况下，PBS的输出完全依赖于光的偏振态。
		
		\item 偏振态的合成
		
		接下来，我们可以探讨是否可以通过混合水平偏振光和垂直偏振光来合成45°偏振光。在理想情况下，如果能够以0.5:0.5的比例混合 \( |H\rangle \) 和 \( |V\rangle \)，理论上可以得到45°偏振光 \( |D\rangle \)。这种合成原理基于量子力学中的叠加态理论。然而，直接混合两个不同方向的偏振光束在实验中往往不太可行，因为光的相位和路径等因素都会对叠加结果产生影响。通常，通过相位调制和偏振旋转等手段更为直接地实现特定偏振态。因此，实验中更常用的方式是通过起偏器和波片来直接制备所需的偏振态。
	\end{itemize}
	
	\item 单光子的探测及相应探测器效率的测量
	
	光子计数是一种用于测量极其微弱光信号的技术，尤其适用于单光子水平的光探测。这种技术具有高度的计数稳定性、较强的抗干扰能力、低噪声、高探测效率等特点，因此被广泛应用于弱光精密测量领域。
	\begin{itemize}
		\item 光电倍增管 (PMT) 的工作原理
		
		单光子的探测通常通过光电倍增管（PMT）等高灵敏度探测器实现。PMT是一种能够将极微弱的光信号（如单个光子）转化为可检测的电信号的器件。其工作原理如下：微弱的光子通过PMT前端的光电窗口入射到光电阴极上，光电阴极受光子的作用会释放出光电子。这些初始光电子经过一系列倍增电极（称为“倍增极”）加速和碰撞，电子的数量被逐级倍增。倍增后的电子最终在阳极处被收集，形成一个电流脉冲。该电流脉冲通过前端放大器放大，随后被转换为一个电压脉冲输出。最终的电压脉冲经过甄别和整形，形成一个固定脉冲幅度和宽度的标准信号，可以用于后续的信号处理和计数。这种信号转换和放大的过程使得PMT能够探测到极其微弱的光信号，甚至单个光子。这种技术被广泛应用于弱光信号的探测场景，如量子光学实验、弱发光检测等。
		
		\begin{figure}[htbp]
			\centering
			\includegraphics[width=0.6\textwidth]{APL1_3_pre_PMT.png}
			\caption{光电倍增管的工作原理}
		\end{figure}
		
		\item 单光子能量的计算
		
		光子的能量 \(E\) 可以通过普朗克常量 \(h\)、光速 \(c\) 和光的波长 \(λ\) 来计算，公式如下：
		\[
		E = h \cdot \nu = \frac{h \cdot c}{\lambda}
		\]
		以404 nm的紫外激光为例，计算单个光子的能量：
		
		\[
		E = \frac{6.626 \times 10^{-34} \cdot 3 \times 10^8}{404 \times 10^{-9}} \approx 4.92 \times 10^{-19} \, \text{J}
		\]
		
		这里，\(h = 6.626 \times 10^{-34} \, \text{J·s}\) 是普朗克常量，\(c = 3 \times 10^8 \, \text{m/s}\) 是光速，\(λ = 404 \, \text{nm}\) 是激光的波长。
		
		\item 平均光子数与功率的计算
		
		在实验中，激光器通常以一定的重复频率发射光脉冲，每个脉冲内包含的平均光子数 \( \mu \) 也可以计算出来。对于一个404 nm的激光源，如果其发射频率为 10 MHz，即每秒发射 \( 10 \times 10^6 \) 个脉冲，当每个脉冲的平均光子数为0.1时，总输出功率 \(P\) 可以通过以下公式计算：
		
		\[
		P = \mu \cdot f \cdot E
		\]
		
		其中，\(f\) 是脉冲激光的发射频率，\(E\) 是每个光子的能量，\(\mu\) 是每个脉冲的平均光子数。将已知数值代入公式：
		
		\[
		P = 0.1 \cdot 10 \times 10^6 \cdot 4.92 \times 10^{-19} = 4.92 \times 10^{-13} \, \text{W}
		\]
		
		由此可以得到，在每个脉冲平均包含0.1个光子的情况下，激光器的输出功率为 \( 4.92 \times 10^{-13} \, \text{W} \)。
		
		\item 衰减计算与单光子制备
		
		实验中，为了制备单光子状态的光源，通常需要通过调整激光器的功率或加入可调的光学衰减器。假设激光器的输出功率为 \( P_0 = 20 \, \mu \text{W} \)，为了将功率衰减至仅包含平均 0.1 个光子/脉冲所需的功率，即 \( P = 4.92 \times 10^{-13} \, \text{W} \)，所需的衰减值可以通过以下公式计算：
		
		\[
		\text{衰减值} = 10 \log_{10} \left( \frac{P_0}{P} \right)
		\]
		
		代入已知数值：
		
		\[
		\text{衰减值} = 10 \log_{10} \left( \frac{20 \times 10^{-6}}{4.92 \times 10^{-13}} \right) \approx 76.1 \, \text{dB}
		\]
		
		这意味着，需要对激光器输出进行 76.1 dB 的衰减，才能得到平均每个脉冲包含0.1个光子的激光输出。通过调节激光器的衰减器，实验人员可以精确控制输出的光子数量。
		
		\item 探测器效率的测量
		
		一旦将激光调节到单光子水平，接下来可以通过单光子探测器来检测这些光子。探测器的计数率 \( N \) 与光子的发射频率 \( f \)、平均光子数 \( \mu \)、以及探测器的效率 \( \eta \) 之间的关系为：
		
		\[
		N = f \cdot \mu \cdot \eta
		\]
		
		在这个公式中，\(f\) 是激光器的发射频率，\( \mu \) 是每个脉冲的平均光子数，\( \eta \) 是探测器的效率。通过实验测量探测器的实际计数 \(N\)，并且知道激光器的发射频率和平均光子数，就可以反推出探测器的效率 \( \eta \)。
		
		例如，若激光器的发射频率为10 MHz，每个脉冲的平均光子数为0.1，探测器的计数率 \( N \) 为 500,000 次/秒，则探测器的效率为：
		
		\[
		\eta = \frac{N}{f \cdot \mu} = \frac{500,000}{10 \times 10^6 \cdot 0.1} = 0.5
		\]
		
		因此，探测器的效率为50\%。
	\end{itemize}
	
	\item 单光子的标定
	
	在单光子量子密钥分发（QKD）实验中，常用的替代方法是采用脉冲激光器，通过强烈的光学衰减，使激光输出的光子数接近单光子水平。虽然这种方法无法完全达到理想单光子源的要求，但在许多实验中被广泛使用，特别是在实现BB84协议的初步实验中。以实验中使用的404nm脉冲激光器为例，其发光频率为10MHz（即每秒发射1000万次脉冲）。通过调节衰减模块（包括固定衰减片和圆形可调衰减片），可以将激光输出的光子数量精确控制到单光子水平。通过调节衰减，激光每个脉冲的平均光子数可以被降低到非常接近1，甚至更低，如0.1或0.05个光子/脉冲。
	
	在理想情况下，单光子光源的标定需要依靠高精度的衰减器以及功率测量装置，例如使用高功率激光器进行初步标定，测量衰减器的实际衰减值。然而，在较为简化的实验中，可以使用单光子探测器来标定衰减器的衰减值，从而简化实验设备的复杂度。
	
	例如，在实验中，如果目标是标定平均每个脉冲包含0.05个光子，并且单光子探测器的探测效率为20\%，则通过探测器的计数率 \(N\) 可以验证衰减是否达到理想水平。具体来说，计数率 \(N\) 可以通过以下公式计算：
	
	\[
	N = f \times \mu \times \eta
	\]
	
	将已知数值代入公式：
	
	\[
	N = 10 \times 10^6 \times 0.05 \times 0.2 = 10^5 \, \text{Hz}
	\]
	
	这意味着，在上述条件下，探测器的计数率应为 \( 10^5 \, \text{Hz} \)（即每秒10万次探测到光子的事件）。通过实际测量探测器的计数率，可以验证衰减器的设置是否达到了期望的单光子水平。
\end{enumerate}


%---------------------------------------------------------------------
% 实验前思考题
\subsection{Thinking Before Experiment}
% 1
\begin{question}
	回顾偏振光实验，说明$\lambda/2$波片，$\lambda/4$波片的工作原理。
\end{question}
相位延迟器是一种光学元件，用于改变穿过它的光波的相位，而不显著改变光波的其他属性，如振幅或频率。这种设备在光学实验、通信以及各种光学仪器中扮演着重要角色，尤其在涉及到控制和操作光的偏振状态的应用中。

相位延迟器通过增加光波不同分量间的相位差来工作。具体来说，它可以改变穿过它的偏振光的相位关系。例如，在一个偏振光波中，可能需要调整垂直分量和平行分量之间的相位差，以实现特定的偏振状态（如将线性偏振光转变为圆偏振光或椭圆偏振光）。

波片是最常见的相位延迟器，它利用双折射材料（如石英或莫尔盐）的性质。这些材料对于垂直方向的光和平行方向的光有不同的折射率，因此不同方向的光波会以不同的速度传播，从而产生相位延迟。波片通常被设计成产生半波长（半波片）或四分之一波长（四分之一波片）的相位差，分别用于实现不同的偏振转换。

\begin{enumerate}
	\item 双折射效应
	
	双折射效应是一种特定类型的光学现象，发生在某些各向异性的材料（如晶体）中。在这些材料中，由于其晶体结构的特性，光在穿过时会根据其偏振方向被分裂为两束具有不同折射率的光。这两束光通常被称为普通光（ordinary ray, o-ray）和非常光（extraordinary ray, e-ray）。
	
	\begin{enumerate}
		\item 普通光（o-ray）：在双折射材料中，普通光的偏振方向垂直于主光轴（晶体的一个特殊方向），它遵循普通的斯涅尔定律，并且具有恒定的折射率，不受光的传播方向影响。
		
		\item 非常光（e-ray）：与普通光不同，非常光的偏振方向平行或部分平行于主光轴。其折射率取决于光的传播方向和主光轴的相对角度，因此不遵循普通的斯涅尔定律。
	\end{enumerate}
	
	双折射效应可以用以下公式进行描述，其中涉及到普通光和非常光的折射率：
	
	\[
	n_o = \text{常数}
	\]
	
	\[
	n_e(\theta) = \frac{n_o n_e}{\sqrt{n_o^2 \sin^2(\theta) + n_e^2 \cos^2(\theta)}}
	\]
	
	其中，\(n_o\) 是普通光的折射率，\(n_e\) 是非常光在主光轴方向的折射率，\(\theta\) 是光的传播方向与主光轴之间的夹角。
	
	当非偏振光进入双折射材料时，光波会根据电场矢量与材料主光轴的相对方向被分解成两个正交偏振分量。这两个分量分别沿着不同的路径传播，因为它们具有不同的折射率，导致相位差的产生，这个相位差可能导致光波在出射时呈现出复杂的偏振状态。		
	
	\item 二分之一波片
	
	二分之一波片（或称半波片）是一种特殊的光学元件，它利用双折射材料的性质来改变通过它的偏振光的偏振状态。通过精确控制光在这种材料中的传播，半波片能够引入一个相对于两个正交偏振分量的相位差，这个相位差正好是半个波长（180度）。这种相位差的引入可以用来旋转偏振光的偏振方向。
	
	假设有一束线性偏振光通过一个半波片，光的电场可以分解为两个分量，这两个分量分别平行和垂直于半波片的光轴。当这两个分量通过半波片时，由于双折射的效果，它们将经历不同的折射率，从而导致一个相位差。
	
	具体来说，半波片引起的相位差 \(\Delta \phi\) 可以表示为：
	
	\[
	\Delta \phi = \frac{2\pi d}{\lambda} (n_e - n_o)
	\]
	
	其中：\(d\) 是半波片的厚度，\(\lambda\) 是光的波长，\(n_e\) 和 \(n_o\) 分别是材料中非常光和普通光的折射率。
	
	为了使得半波片达到其功能，我们需要调整其厚度 \(d\)，使得 \(\Delta \phi = \pi\)（相当于180度），这对应于光波的半个波长。当相位差为180度时，通过半波片的线性偏振光的偏振方向将会旋转90度。
	
	\begin{figure}[htbp]
		\centering
		\includegraphics[width=0.4\textwidth]{Lab2_6Gra13.jpg}
		\caption{二分之一波片\cite{c}}
	\end{figure}
	
	设原始偏振光的电场向量为 \(\vec{E}_0\)，假设它与半波片的一个光轴（例如，快轴）成角 \(\theta\)。通过半波片后，偏振方向将旋转至原始方向的正交方向。数学上，可以这样描述这一过程：
	
	假设初始电场向量为：
	\[
	\vec{E}_0 = E_0 \cos(\theta) \hat{i} + E_0 \sin(\theta) \hat{j}
	\]
	
	经过半波片后，\(\hat{i}\) 方向的分量经历相位延迟，变为：
	\[
	E_0 \cos(\theta) \hat{i} e^{i\pi}
	\]
	
	由于 \(e^{i\pi} = -1\)，输出的电场向量将是：
	\[
	\vec{E} = -E_0 \cos(\theta) \hat{i} + E_0 \sin(\theta) \hat{j}
	\]
	
	这意味着原来沿 \(\hat{i}\) 方向的分量反向，而 \(\hat{j}\) 方向的分量保持不变，从而导致偏振方向旋转了90度。
	
	
	\item 四分之一波片
	
	四分之一波片（或称四分之一波长片，简称$\lambda/4$片）是一种特殊的光学元件，用于将线性偏振光转换成圆偏振光，或反之。与半波片类似，四分之一波片利用双折射材料的性质来引入一个相对于两个正交偏振分量的相位差，不过在这种情况下，相位差正好是四分之一波长（90度）。
	
	四分之一波片的工作原理基于其能够为垂直和平行于其光轴的偏振光分量间引入90度的相位差。这种相位差使得线性偏振光转化为圆偏振光，或将圆偏振光转化为线性偏振光，具体取决于光的初始状态和波片的方向。
	
	相位差 \(\Delta \phi\) 引入公式为：
	\[
	\Delta \phi = \frac{2\pi d}{\lambda} (n_e - n_o)
	\]
	其中：\(d\) 是波片的厚度，\(\lambda\) 是光的波长，\(n_e\) 和 \(n_o\) 分别是材料中非常光和普通光的折射率。
	
	为了使得四分之一波片起作用，需要调整其厚度 \(d\) 以确保 \(\Delta \phi = \frac{\pi}{2}\) （相当于90度）。这会引起两个正交偏振分量间四分之一波长的相位延迟。
	
	\begin{figure}[htbp]
		\centering
		\includegraphics[width=0.4\textwidth]{Lab2_6Gra14.jpg}
		\caption{四分之一波片\cite{c}}
	\end{figure}
	
	假设初始的线性偏振光沿特定方向进入四分之一波片，可以假设电场矢量 \(\vec{E}_0\) 成角 \(\theta\) 与快轴。通过波片后，其中一个分量（假设是快轴对应的分量）会比另一个分量（慢轴对应的分量）先行一个四分之一波长的相位。
	
	电场分量可以表示为：
	
	\begin{align*}
		E_x &= E_0 \cos(\theta) e^{i\omega t} \\
		E_y &= E_0 \sin(\theta) e^{i(\omega t + \pi/2)}
	\end{align*}
	
	其中 \(E_x\) 和 \(E_y\) 分别是沿快轴和慢轴的电场分量。
	
	由于引入了\(\pi/2\)的相位差，这将导致电场矢量旋转，形成圆偏振光。如果初始偏振与快轴和慢轴对齐，则可以得到右旋或左旋圆偏振光，具体取决于相位的领先或落后关系。
	
\end{enumerate}

% 2
\begin{question}
	如何检测一个任意方向的线偏振光？
\end{question}
\begin{enumerate}
	\item 使用起偏器 (Polarizer)
	
	起偏器是一种只能透射某一方向偏振光的光学器件。通过起偏器的角度调整，可以检测到入射光的偏振状态。具体方法如下：将待测光束通过一个可旋转的起偏器。起偏器只允许某一方向的线偏振光通过，因此可以检测到入射光的不同偏振分量。当旋转起偏器时，可以用光探测器测量通过的光强度。根据强度的变化，可以判断入射光的偏振角。
	
	\begin{figure}[htbp]
		\centering
		\includegraphics[width=0.5\textwidth]{Lab2_6Gra4.jpg}
		\caption{自然光入射到线起偏器\cite{c}}
		\label{fig:fig4}
	\end{figure}
	
	根据马吕斯定律 (Malus’s Law)，通过起偏器后透射光的强度 \( I \) 与起偏器角度 \( \theta \)（相对于光的偏振方向）之间的关系为：
	\[
	I = I_0 \cos^2 \theta
	\]
	其中 \( I_0 \) 是光的初始强度，\( \theta \) 是入射光偏振方向与起偏器传输轴之间的夹角。通过转动起偏器并测量光强度，便可以根据强度的最大值（当起偏器的传输轴与偏振方向一致时）来确定偏振光的方向。
	
	\begin{figure}[htbp]
		\centering
		\includegraphics[width=0.6\textwidth]{Lab2_6Gra6.jpg}
		\caption{Malus定律\cite{c}}
	\end{figure}
	
	\item 偏振分光棱镜 (PBS)
	
	偏振分光棱镜 (PBS) 是一种可以将入射光分成两个正交偏振方向（通常是水平偏振和垂直偏振）的光学器件。具体检测步骤如下：入射光通过PBS：将待测光束垂直入射到PBS上，PBS会将光束分成两部分：水平偏振光 \( |H\rangle \) 和垂直偏振光 \( |V\rangle \)。探测分光后的光强：通过探测器分别测量透射光和反射光的强度，即水平偏振光和垂直偏振光的分量。
	
	\begin{figure}[htbp]
		\centering
		\includegraphics[width=0.35\textwidth]{Lab2_6Gra11.jpg}
		\caption{Wollaston棱镜\cite{c}}
	\end{figure}
	
	对于任意方向的线偏振光，可以表示为水平和垂直偏振光的叠加态：
	\[
	|E\rangle = E_H |H\rangle + E_V |V\rangle
	\]
	其中 \( E_H \) 和 \( E_V \) 分别是水平偏振和垂直偏振的振幅。通过测量PBS透射和反射光的强度比，可以推导出入射光的偏振方向。例如，如果透射光和反射光的强度相等，则说明入射光是45°的线偏振光。
	
	\item 使用半波片和PBS组合
	
	使用半波片和PBS的组合可以更精确地测量偏振角度。步骤如下：半波片调节偏振方向：半波片是一种光学器件，它可以通过旋转调节入射光的偏振方向。当光通过旋转角度为 \( \theta \) 的半波片时，其偏振方向会发生 \( 2\theta \) 的旋转。将任意方向的线偏振光经过半波片，调整半波片的角度，可以将偏振光旋转至与PBS的分光方向一致（水平或垂直方向）。PBS分光检测：将经过半波片的光束送入PBS中，通过测量PBS的透射光和反射光，可以确定入射光的偏振方向。例如，如果半波片旋转了一个角度，使得光通过PBS后只有一个方向有输出（透射或反射），则可以判断出入射光的偏振方向与半波片的旋转角度有直接关系。
	
	\item 斯托克斯参数法
	
	斯托克斯参数法是一种更加系统的偏振态测量方法。斯托克斯参数包括四个值 \( S_0 \)、\( S_1 \)、\( S_2 \)、\( S_3 \)，它们可以完全描述任意偏振光的状态。具体来说，斯托克斯参数通过多个角度的偏振测量得到。
	
	\( S_0 \)：总强度，即不受偏振影响的光强度。\( S_1 \)：水平偏振光与垂直偏振光的强度差异。\( S_2 \)：45°偏振光与135°偏振光的强度差异。\( S_3 \)：右旋圆偏振光与左旋圆偏振光的强度差异。
	
	为了检测任意方向的线偏振光，需要依次测量水平偏振光、垂直偏振光、45°偏振光和135°偏振光的强度，并通过这些测量结果计算出斯托克斯参数。具体步骤为：首先，通过PBS或起偏器测量水平和垂直偏振光的强度，得到 \( S_1 \)。接着，通过旋转45°的起偏器测量45°和135°偏振光的强度，得到 \( S_2 \)。最后，如果需要区分圆偏振光，可以进一步测量右旋和左旋圆偏振光，得到 \( S_3 \)。
\end{enumerate}

% 3
\begin{question}
	单光子为什么不能直接用普通功率计测量？
\end{question}
单光子不能直接用普通功率计测量的原因在于单光子信号的极低能量和功率，超出了普通功率计的测量能力。

单光子的能量由普朗克常数 \( h \) 和光子的频率 \( \nu \) 决定，公式为：
\[
E = h \cdot \nu = \frac{h \cdot c}{\lambda}
\]
其中：
- \( h = 6.626 \times 10^{-34} \, \text{J·s} \) 是普朗克常数，
- \( c = 3 \times 10^8 \, \text{m/s} \) 是光速，
- \( \lambda \) 是光的波长。

例如，假设我们使用波长为404 nm的紫外光，则单个光子的能量为：
\[
E = \frac{6.626 \times 10^{-34} \cdot 3 \times 10^8}{404 \times 10^{-9}} \approx 4.92 \times 10^{-19} \, \text{J}
\]
这一能量极其微小，仅为$10^-19$ J的量级，这比普通功率计的灵敏度要低得多。普通功率计通常用于测量毫瓦或微瓦级的光功率，而单光子的能量要比这低出许多数量级。

单光子对应的极低功率，如果我们以一个典型的激光脉冲重复频率来计算单光子对应的功率，假设激光器的脉冲频率为10 MHz（即每秒产生 \( 10^7 \) 个脉冲），并且每个脉冲内只有1个光子，单光子的平均功率为：
\[
P = E \cdot f = 4.92 \times 10^{-19} \, \text{J} \times 10^7 \, \text{Hz} = 4.92 \times 10^{-12} \, \text{W} = 4.92 \, \text{pW}
\]
这个功率值为皮瓦（pW）级别，而普通功率计通常只能检测到微瓦（\(\mu W\)）甚至毫瓦（mW）级别的功率。因此，单光子的功率远低于普通功率计的探测极限。

普通光功率计的工作原理基于吸收或转换光能量为热能或电能，进而测量这些能量的变化。具体来说，普通功率计通常使用热电堆或光电二极管来探测光的功率：基于热效应，光能量通过吸收材料转化为热量，热电材料产生电压与功率成正比。热电堆对能量的变化有一定的响应时间，通常用于测量连续光源的平均功率，无法对极低能量的单光子进行快速响应。通过光电效应将光信号转换为电信号，电流的大小与入射光的功率成正比。然而，光电二极管的探测灵敏度受到噪声和响应时间的限制，通常无法检测到单个光子的极低功率。由于单光子的能量远低于普通功率计的探测下限，普通功率计无法产生足够的信号响应来检测单光子。

与普通功率计相比，单光子探测器（Single-Photon Detector, SPD）具有更高的灵敏度，专门设计用于探测极低强度的光信号。这类探测器通常基于光电倍增管（PMT）或雪崩光电二极管（APD）等技术。光电倍增管 (PMT)：PMT能够探测单个光子并将其放大为可测的电信号。光子入射到光电阴极时，会释放出电子，随后通过多个倍增极逐步放大，最终被阳极收集成可检测的电流脉冲。这种高灵敏度的放大机制使PMT能够探测单光子事件。雪崩光电二极管 (APD)：APD在工作于盖革模式（Geiger mode）时，能够通过光电效应触发雪崩倍增，放大单光子的电子信号，产生明显的电流脉冲。APD具有极高的光子探测效率，适合用于单光子计数和量子光学实验。这些单光子探测器的设计目标就是为了在极低光功率条件下产生明显的电信号，因此它们能够在单光子水平下工作，而普通功率计无法做到这一点。

在经典光学中，光可以被看作是连续的电磁波，功率可以表示为光束中单位时间内传递的能量。然而，在量子力学框架下，光表现为一系列离散的光子，每个光子的能量都极小。普通功率计是为检测宏观光功率设计的，它基于经典光学原理，无法感知单光子的量子特性。相反，单光子探测器则基于光子的量子效应进行探测。它们能够感知每个单独的光子事件，并且以脉冲的形式输出信号。这种差异使得单光子无法直接通过普通功率计来测量。

% 4
\begin{question}
	检验单光子探测器的探测效率可以用强光吗？
\end{question}
检验单光子探测器的探测效率不能直接使用强光，因为单光子探测器是专门设计用于探测极低光强度的单光子事件。强光会导致探测器进入过饱和状态，甚至可能损坏探测器。

单光子探测器的设计目标是探测单个光子事件，它们通过高度敏感的探测机制，如光电倍增管 (PMT) 或雪崩光电二极管 (APD)，将单个光子的电磁信号放大为可检测的电信号。这类探测器通常在光强极低的条件下工作，在单光子级别的光强下，每个光子被探测后产生一个脉冲信号。这些探测器专为单光子量级的光信号设计，具有高灵敏度。但在强光照射下，探测器不仅无法精确区分单个光子，还可能被过载或损坏。

单光子探测器通常设计用于处理极低的光子输入。如果用强光照射探测器，会产生大量光子，这会导致探测器超出其工作范围，并进入“饱和”状态。在这种状态下，探测器将无法区分单个光子事件，探测信号会变得饱和，输出的脉冲信号失去与输入光子数量的线性关系。饱和效应通常表现为探测器的输出信号不再增加，甚至不再响应输入光子。这意味着探测器无法准确测量光信号，无法评估其探测效率。

许多单光子探测器对光输入非常敏感，强光可能会导致过大的电流流过探测器的内部组件，导致探测器损坏。特别是雪崩光电二极管（APD）在强光下可能会产生过大的雪崩效应，长时间暴露于强光下会导致器件永久性损坏。

即使强光不会完全使探测器饱和，它也会增加背景噪声。强光会触发大量电子，产生高水平的暗噪声或热噪声，从而掩盖单光子信号。这会导致探测器的信噪比显著降低，增加探测效率测量的误差。

% 5
\begin{question}
	BB84 协议的原理和步骤。
\end{question}
在\textbf{原理概述}部分有详细的叙述，以下是总结：
\begin{itemize}
	\item BB84协议的安全性基于以下量子力学原理：
	
	量子态不可克隆定理：未知的量子态无法被精确复制。这意味着潜在的窃听者无法无痕复制光子的量子态，而不引起原始量子态的扰动。
	
	测量引起的扰动：对量子态的测量会不可避免地改变其状态。如果窃听者尝试测量光子的偏振态，会引入可检测的错误。
	
	量子叠加和正交性：不同基中的量子态是不可区分的。例如，在两个不同的偏振基（如矩形基和对角基）中，相应的态是正交的，无法在不知基的情况下准确测量。
	
	\item BB84协议主要包括以下步骤，通常在两个通信方之间进行，称为Alice（发送方）和Bob（接收方）：
	\begin{enumerate}
		\item 密钥编码
		
		\begin{itemize}
			\item 选择基和比特值：
			\begin{itemize}
				\item Alice 随机选择一系列比特值（0和1）。
				\item 对于每个比特，Alice 随机选择一个偏振基进行编码，通常是两个互相不兼容的基：
				\begin{itemize}
					\item 矩形基（+基）：水平（0°）和垂直（90°）偏振，分别表示比特0和1。
					\item 对角基（×基）：45°和135°偏振，分别表示比特0和1。
				\end{itemize}
			\end{itemize}
			
			\item 量子态发送：Alice 根据选择的基和比特值，生成相应偏振态的光子序列，逐个发送给 Bob。
		\end{itemize}
		
		\item 测量和基选择
		
		\begin{itemize}
			\item Bob 的基选择：
			\begin{itemize}
				\item Bob 在接收到每个光子时，随机选择一个偏振基（+基或×基）进行测量。
				\item Bob 记录每个测量的结果和所使用的基。
			\end{itemize}
		\end{itemize}
		
		\item 基的公开和筛选
		
		\begin{itemize}
			\item 基信息的交换：
			\begin{itemize}
				\item 传输完成后，Alice 和 Bob 通过公共信道（可以被窃听但无法篡改）公开各自使用的基序列，但不公开比特值。
			\end{itemize}
			
			\item 基的匹配：
			\begin{itemize}
				\item Alice 和 Bob 比较各自的基序列，保留基匹配的位置，丢弃不匹配的位置。
				\item 由于基的随机选择，理论上有50%的位置基匹配。
			\end{itemize}
		\end{itemize}
		
		\item 初始密钥的生成
		
		\begin{itemize}
			\item 提取初始密钥：
			\begin{itemize}
				\item 在基匹配的位置，Bob 保留测量结果，形成一个初始密钥序列。
				\item Alice 也根据相同的位置提取相应的比特值。
			\end{itemize}
		\end{itemize}
		
		\item 误码率检测
		
		\begin{itemize}
			\item 抽样比较：
			\begin{itemize}
				\item Alice 和 Bob 从初始密钥中随机抽取一部分比特，公开并比较这些比特值。
				\item 通过比较，计算量子误码率（QBER），即比特错误的比例。
			\end{itemize}
			
			\item 安全性评估：
			\begin{itemize}
				\item 如果 QBER 低于某个阈值（通常为11%），则认为密钥基本安全，可以继续。
				\item 如果 QBER 过高，说明可能存在窃听或信道噪声太大，需要舍弃密钥并重新开始。
			\end{itemize}
		\end{itemize}
		
		\item 信息协调
		
		错误纠正（Error Correction）：使用经典的纠错协议（如 Cascade、BCH 码等），Alice 和 Bob 通过公共信道协商，纠正初始密钥中的错误比特，确保双方密钥一致。
		
		
		\item 隐私放大
		
		\begin{itemize}
			\item 抵御信息泄露：
			\begin{itemize}
				\item 由于可能存在窃听者（Eve）获取了部分信息，Alice 和 Bob 需要进行隐私放大。
				\item 通过应用哈希函数等方法，将初始密钥压缩为一个更短的最终密钥，减少 Eve 已知信息的影响。
			\end{itemize}
		\end{itemize}
		
		\item 最终密钥的生成
		
		安全密钥：经过上述步骤，Alice 和 Bob 获得了一致且保密的最终密钥，可以用于加密通信（如一次一密）。
		
	\end{enumerate}
	
\end{itemize}


% 6
\begin{question}
	密钥分发过程中，为什么需要有同步信号？
\end{question}
\begin{itemize}
	\item 精确时间对齐
	
	在QKD协议（如BB84协议）中，Alice（发送方）和Bob（接收方）需要以非常高的时间精度来处理光子传输和检测。同步信号的主要作用是确保Alice发出的每个光子和Bob接收到的每个光子能够在时间上正确匹配。光子通过光纤或自由空间传输的速度非常快（接近光速），传输时间极短。如果没有时间同步，Bob很难确定他在何时应该开始探测光子。在实际QKD系统中，Bob通常会在特定的时间窗口内探测光子，如果光子到达的时间超出这个窗口期，Bob将无法准确地探测到光子。例如，光子可能会因为信道噪声、光纤长度差异等原因稍有延迟。如果Alice和Bob没有同步时间，Bob可能会错过某些光子或者将它们归到错误的时间点。探测器在无光子到达的情况下也可能会产生“暗计数”噪声。如果没有时间同步，Bob可能会将噪声事件误认为是Alice发送的光子，从而增加误码率（QBER）。
	
	\item 数据关联与基选择对齐
	
	在QKD过程中，Alice会发送一系列量子态（例如偏振态），这些量子态对应于她随机生成的比特序列。Bob在接收端会随机选择测量基进行探测，并记录下测量结果。在密钥提取过程中，Alice和Bob需要将发送的量子态与接收的测量结果进行一一对应，这就需要同步信号来帮助他们进行数据关联。在BB84协议中，Alice和Bob需要在公共信道上公开各自选择的偏振基。为了确保双方在同一个时间点选择了相同的基，他们必须知道在某个时刻Alice发出的光子对应的是Bob在何时测量到的光子。如果没有时间同步，Alice和Bob的基选择可能会错位，导致本应对应的比特无法匹配，从而增加误码率，影响最终密钥的长度和质量。
	
	\item 高效数据处理
	
	QKD系统涉及大量的光子传输和探测，光子是以非常高的频率（通常是MHz级别）发射和接收的。没有时间同步，Bob在接收端将很难将每次接收到的光子与Alice发出的光子进行快速匹配。为了有效管理大量的数据流，Alice和Bob需要时间同步来确保数据处理的高效性。如果没有时间同步，Bob可能需要花费大量的时间和计算资源来进行数据匹配，降低系统的效率。Bob的探测器和电子电路需要与Alice的光子发射保持同步。如果不同步，Bob的探测窗口将与Alice的光子发射不同步，导致数据丢失或产生错误数据。
	
	\item 降低误码率和提高密钥生成率
	
	时间同步在QKD中对减少误码率（QBER）和提高密钥生成率至关重要。原因如下：
	
	如果Alice和Bob在时间上不同步，Bob可能会错误地测量光子的到达时间，从而影响测量结果。这种情况下，本应是有效的光子信息可能会被误认为是噪声或误码。QKD中的误码率越高，最终用于加密通信的密钥生成率就越低。较高的误码率不仅影响密钥的安全性，还会增加后续纠错和隐私放大步骤的计算负担。
	
	通过精确的时间同步，Bob可以在正确的时间窗口内探测光子，减少探测错误。这样可以确保更多的有效光子被正确接收，增加最终密钥生成的效率。如果没有时间同步，很多有效光子可能会因为Bob的探测窗口偏移而未被检测到，导致光子浪费。时间同步能够确保Alice发送的每个光子尽可能地被Bob接收到。
	
	\item 防止窃听攻击
	
	时间同步在防止窃听攻击方面也起到了关键作用。量子密钥分发的安全性依赖于窃听者无法无痕复制和拦截光子。如果Alice和Bob没有精确的时间同步，窃听者可能利用这一点进行攻击，具体如下：
	
	窃听者（Eve）可能尝试通过引入时延或伪造时间戳来破坏Alice和Bob的同步，进而破坏他们的密钥生成过程。如果Alice和Bob不同步，他们可能无法检测到Eve的存在，因为错误可能会被解释为时间误差而不是窃听行为。
	
	如果时间同步做得非常精确，Alice和Bob可以通过检测比特错误率（QBER）的变化来判断是否存在窃听者。不同步的时间信号可能会增加错误率，而窃听者引入的扰动也会增加QBER。如果没有良好的同步，Alice和Bob无法区分这些错误是由窃听者引起的，还是由时间偏差引起的。
	
	\item 量子信号与经典信号的配合
	
	在QKD系统中，量子信号（即光子传输）和经典信号（基选择、误码率比较、纠错等）需要相互配合。Alice发送光子，Bob需要在正确的时间窗口内探测光子。Alice和Bob需要通过经典信道交换基信息、错误校正信息等。这些信息的交换需要与量子信号的传输时间保持一致，以确保Alice和Bob在相同的时间序列上操作。如果量子信号和经典信号不同步，Alice和Bob将无法正确关联每个光子事件，导致密钥无法正确生成。
\end{itemize}

